% =====================================================================
% fig_authoring_pipeline.tex -- how a physics description in text becomes the
% control video the diffusion backbone follows: a language model writes the
% scene specification, the specification is simulated and validated on the
% trajectory it actually produces, a failing specification is returned to the
% model with the violated condition attached and re-authored, and only an
% accepted specification is rendered.
%
% \input-able into main.tex, which loads adobe-techreport.sty (hence the
% adobered / adobeink / adobegray colours and the Source font family).
%
% REQUIRED IN THE MAIN PREAMBLE:
%     \usepackage{tikz}
%     \usetikzlibrary{arrows.meta,positioning,calc}
% Both lines are required, and the libraries MUST be loaded in the preamble --
% see the guard below for why a figure file cannot load them for itself.
% Every style is declared inside the tikzpicture's own option list, so all
% definitions are local to this picture: \input-ing this file twice, or
% alongside its companion figures, cannot clash.
%
% CONVENTIONS shared with fig_critic_pipeline.tex and fig_graph_construction.tex:
%   adobegray!12 fill      = a learned component
%   white + adobeink line  = a deterministic component
%   adobered line          = the one check the pipeline turns on (used here for
%                            validation of the simulated trajectory, and for
%                            nothing else)
%   dashed gray            = something handed to the system, not produced by it
% Every stage in this picture is implemented, so every line is solid; the
% solid/dashed built-versus-roadmap distinction used elsewhere in the report
% does not apply here.
%
% WRAPPING: like every other figs/*.tex in this report, this file contains no
% figure environment and no \caption -- the caller supplies both:
%     \begin{figure}[tp]
%       \centering
%       % =====================================================================
% fig_authoring_pipeline.tex -- how a physics description in text becomes the
% control video the diffusion backbone follows: a language model writes the
% scene specification, the specification is simulated and validated on the
% trajectory it actually produces, a failing specification is returned to the
% model with the violated condition attached and re-authored, and only an
% accepted specification is rendered.
%
% \input-able into main.tex, which loads adobe-techreport.sty (hence the
% adobered / adobeink / adobegray colours and the Source font family).
%
% REQUIRED IN THE MAIN PREAMBLE:
%     \usepackage{tikz}
%     \usetikzlibrary{arrows.meta,positioning,calc}
% Both lines are required, and the libraries MUST be loaded in the preamble --
% see the guard below for why a figure file cannot load them for itself.
% Every style is declared inside the tikzpicture's own option list, so all
% definitions are local to this picture: \input-ing this file twice, or
% alongside its companion figures, cannot clash.
%
% CONVENTIONS shared with fig_critic_pipeline.tex and fig_graph_construction.tex:
%   adobegray!12 fill      = a learned component
%   white + adobeink line  = a deterministic component
%   adobered line          = the one check the pipeline turns on (used here for
%                            validation of the simulated trajectory, and for
%                            nothing else)
%   dashed gray            = something handed to the system, not produced by it
% Every stage in this picture is implemented, so every line is solid; the
% solid/dashed built-versus-roadmap distinction used elsewhere in the report
% does not apply here.
%
% WRAPPING: like every other figs/*.tex in this report, this file contains no
% figure environment and no \caption -- the caller supplies both:
%     \begin{figure}[tp]
%       \centering
%       % =====================================================================
% fig_authoring_pipeline.tex -- how a physics description in text becomes the
% control video the diffusion backbone follows: a language model writes the
% scene specification, the specification is simulated and validated on the
% trajectory it actually produces, a failing specification is returned to the
% model with the violated condition attached and re-authored, and only an
% accepted specification is rendered.
%
% \input-able into main.tex, which loads adobe-techreport.sty (hence the
% adobered / adobeink / adobegray colours and the Source font family).
%
% REQUIRED IN THE MAIN PREAMBLE:
%     \usepackage{tikz}
%     \usetikzlibrary{arrows.meta,positioning,calc}
% Both lines are required, and the libraries MUST be loaded in the preamble --
% see the guard below for why a figure file cannot load them for itself.
% Every style is declared inside the tikzpicture's own option list, so all
% definitions are local to this picture: \input-ing this file twice, or
% alongside its companion figures, cannot clash.
%
% CONVENTIONS shared with fig_critic_pipeline.tex and fig_graph_construction.tex:
%   adobegray!12 fill      = a learned component
%   white + adobeink line  = a deterministic component
%   adobered line          = the one check the pipeline turns on (used here for
%                            validation of the simulated trajectory, and for
%                            nothing else)
%   dashed gray            = something handed to the system, not produced by it
% Every stage in this picture is implemented, so every line is solid; the
% solid/dashed built-versus-roadmap distinction used elsewhere in the report
% does not apply here.
%
% WRAPPING: like every other figs/*.tex in this report, this file contains no
% figure environment and no \caption -- the caller supplies both:
%     \begin{figure}[tp]
%       \centering
%       % =====================================================================
% fig_authoring_pipeline.tex -- how a physics description in text becomes the
% control video the diffusion backbone follows: a language model writes the
% scene specification, the specification is simulated and validated on the
% trajectory it actually produces, a failing specification is returned to the
% model with the violated condition attached and re-authored, and only an
% accepted specification is rendered.
%
% \input-able into main.tex, which loads adobe-techreport.sty (hence the
% adobered / adobeink / adobegray colours and the Source font family).
%
% REQUIRED IN THE MAIN PREAMBLE:
%     \usepackage{tikz}
%     \usetikzlibrary{arrows.meta,positioning,calc}
% Both lines are required, and the libraries MUST be loaded in the preamble --
% see the guard below for why a figure file cannot load them for itself.
% Every style is declared inside the tikzpicture's own option list, so all
% definitions are local to this picture: \input-ing this file twice, or
% alongside its companion figures, cannot clash.
%
% CONVENTIONS shared with fig_critic_pipeline.tex and fig_graph_construction.tex:
%   adobegray!12 fill      = a learned component
%   white + adobeink line  = a deterministic component
%   adobered line          = the one check the pipeline turns on (used here for
%                            validation of the simulated trajectory, and for
%                            nothing else)
%   dashed gray            = something handed to the system, not produced by it
% Every stage in this picture is implemented, so every line is solid; the
% solid/dashed built-versus-roadmap distinction used elsewhere in the report
% does not apply here.
%
% WRAPPING: like every other figs/*.tex in this report, this file contains no
% figure environment and no \caption -- the caller supplies both:
%     \begin{figure}[tp]
%       \centering
%       \input{figs/fig_authoring_pipeline}
%       \caption{...}
%     \end{figure}
% The \label is carried HERE (see the block just below \makeatother), so main.tex
% must NOT add a second \label{fig:authoring}.
%
% CONTENT follows the subsection "From a physics description to a control video"
% (\ref{sec:gen-authoring}) and is drawn from the implementation in
% ~/diffphy_exp014/scripts/exp014/agentic_mask.py -- stage order from build_one()
% at lines 589-663, with:
%   text-in, no reference video   docstring l.4-8; author() l.294-306; main() l.683
%   scene specification written   author() l.294-327; scene contract l.124-149
%   required fields re-asked once _validate() l.249-265; author() l.322-326
%   waypoints only where needed   docstring l.26-29; motion model l.98-117
%   simulated by the renderer's   _sim_traj() l.348-352 (expand + simulate);
%     own engine                  silhouette_render.render_frames l.209 calls the
%                                 same simulate(); physics_world.py is pymunk
%   validation on the trajectory  lint_scene() l.355-450 (ballistic must be
%                                 dynamic l.406; off-frame l.410-417; rest /
%                                 support l.419-424; rising must rise l.426-428;
%                                 reach / miss targets l.430-445)
%   repair loop, three rounds     build_one() l.604-612; repair() l.472-488;
%                                 realized-motion readout _traj_summary l.453-469
%   appearance written apart      author_texture() l.272-291; called l.616
%   camera + graphic-body strip   apply_camera() l.491-504; _strip_graphic_bodies
%                                 l.210-215; called l.618-619
%   render + contact sheet        build_one l.623-627; render_gif; strip l.507-521
%   what the generator receives   build_one l.632-662 (fill-ratio cap l.632-638,
%                                 low-control-strength cap l.652-654), consumed by
%                                 generate_worker.py l.176-183 -> video_stylizer
%                                 WanVACEStylizer.stylize l.58-135
% The figure is descriptive of the method only: no accuracy, no comparison
% between configurations, and no quantity of any kind appears in it.
%
% The picture is a horizontal left-to-right pipeline about 12.9cm wide, so it
% stays inside the report's text width unscaled -- do not add scale= or
% \resizebox. Box labels are intentionally terse: the operational detail lives in
% the prose of \ref{sec:gen-authoring}, \ref{sec:gen-control} and
% \ref{sec:gen-mechanism}, so the figure carries only the dataflow -- the main
% spine (author -> simulate -> validate -> render -> generate), the bounded
% re-author loop dropping below the simulate/validate span, and the separate
% appearance track feeding the hand-off from below.
% =====================================================================
% The TikZ libraries must be loaded in the PREAMBLE, not here: a figure
% environment forms a group, and TikZ defines a library's macros locally, so a
% \usetikzlibrary issued inside a figure does not survive to the next one. This
% guard turns a missing preamble line into a clear message instead of TikZ's
% confusing "You need to say \usetikzlibrary{calc}" further down.
\makeatletter
\@ifundefined{tikz@library@calc@loaded}{%
  \PackageError{fig_authoring_pipeline}{%
    The preamble must contain\MessageBreak
    \string\usetikzlibrary{arrows.meta,positioning,calc}%
  }{Add that line after \string\usepackage{tikz}.}}{}
\makeatother

% The caller's \input runs BEFORE its \caption, so a bare \label here would bind
% to the enclosing section instead of to the figure. Step the figure counter,
% label that value, then hand the counter back: the \caption that follows steps
% it again to the same number, so \cref{fig:authoring} resolves to this figure.
% Outside a float the block does nothing.
\makeatletter
\@ifundefined{@captype}{}{%
  \refstepcounter{figure}%
  \label{fig:authoring}%
  \addtocounter{figure}{-1}%
}
\makeatother

\begingroup
\providecommand{\apsub}[1]{{\scriptsize\color{adobegray}#1}}

\begin{tikzpicture}[
    x=1cm, y=1cm,
    % --- shared box shape: terse title only; detail lives in the prose ----
    apbox/.style   = {rounded corners=2pt, line width=0.7pt, align=center,
                      inner xsep=6pt, inner ysep=4.2pt,
                      font=\sffamily\footnotesize, text=adobeink,
                      minimum height=8.5mm},
    % --- LEARNED component: light gray fill --------------------------------
    aplearn/.style = {apbox, draw=adobeink, fill=adobegray!12},
    % --- DETERMINISTIC component: white with an ink outline ----------------
    apdet/.style   = {apbox, draw=adobeink, fill=white},
    % --- accent: the check a trajectory must pass before it is rendered -----
    apgate/.style  = {apbox, draw=adobered, line width=0.9pt, fill=white,
                      text=adobered},
    % --- what is handed to the system (input) ------------------------------
    apdata/.style  = {apbox, draw=adobegray, fill=white,
                      dash pattern=on 2pt off 1.6pt},
    apline/.style  = {draw=adobeink, line width=0.7pt},
    aparr/.style   = {apline, -{Straight Barb[length=3.6pt,width=3.8pt]}},
    % the re-author loop, drawn in the accent colour so the cycle reads at a glance
    aploop/.style  = {draw=adobered, line width=0.7pt,
                      -{Straight Barb[length=3.6pt,width=3.8pt]}},
    apnote/.style  = {font=\sffamily\scriptsize, text=adobegray,
                      inner sep=0pt, align=left},
    aptag/.style   = {font=\sffamily\scriptsize, text=adobegray, fill=white,
                      inner xsep=3pt, inner ysep=1pt, align=center},
    aptagr/.style  = {aptag, text=adobered},
    apsw/.style    = {rounded corners=1.2pt, line width=0.7pt, inner sep=0pt,
                      minimum width=11pt, minimum height=7.5pt},
  ]

  % Layout: a horizontal main spine along y=0, left to right, sized to the text
  % width (five boxes: author -> simulate -> validate -> render -> generate).
  % The hand-off is folded into the render->generate arrow. The bounded
  % re-author loop drops below the simulate/validate span; the appearance track
  % sits below the render->generate span and joins the hand-off arrow from below.
  \def\bw{1.8cm}          % spine-box width
  \def\hg{0.46}           % horizontal gap between spine boxes
  \tikzset{spine/.style={apbox, text width=\bw, minimum height=13mm,
                         anchor=west}}

  % ---------------- main spine (left -> right) -----------------------------
  \node[spine, apdata] (spec)   at (0,0) {\textbf{Physics request}\\[1pt]\apsub{text}};
  \node[spine, aplearn] (author) at ($(spec.east)+(\hg,0)$)
    {\textbf{LM writes scene spec}};
  \node[spine, apdet] (sim) at ($(author.east)+(\hg,0)$)
    {\textbf{Simulate}\\[1pt]\apsub{renderer's engine}};
  \node[spine, apgate] (gate) at ($(sim.east)+(\hg,0)$)
    {\textbf{Validate trajectory}};
  \node[spine, apdet] (render) at ($(gate.east)+(\hg,0)$)
    {\textbf{Camera + render control}};
  \node[spine, aplearn] (gen) at ($(render.east)+(1.35,0)$)
    {\textbf{Diffusion backbone}};

  % ---------------- re-author loop (accent), below sim/validate ------------
  \node[aplearn, text width=3.4cm, anchor=north, minimum height=8.5mm]
    (repair) at ($(sim.south)!0.5!(gate.south)+(0,-0.92)$)
    {\textbf{Re-author}\ \ \apsub{$\leq 3$ rounds}};

  % ---------------- appearance track (learned), below render->gen ----------
  \node[aplearn, text width=3.7cm, anchor=north, minimum height=8.5mm]
    (texture) at ($(render.south east)!0.5!(gen.south west)+(0,-0.92)$)
    {\textbf{Appearance written separately}\\[1pt]\apsub{no access to the control}};

  % ---------------- edges: main spine --------------------------------------
  \draw[aparr] (spec)    -- (author);
  \draw[aparr] (author)  -- (sim);
  \draw[aparr] (sim)     -- (gate);
  \draw[aparr] (gate)    -- (render)
        node[aptagr, midway, above=0pt] {validated};
  % render -> generate carries the hand-off (control + appearance + seed + strength)
  \draw[aparr] (render)  -- (gen)
        node[aptag, midway, above=0pt] {hand off};

  % ---------------- re-author loop: validate -> re-author -> re-simulate ----
  \draw[aploop] (gate.south) |- (repair.east)
        node[aptagr, pos=0.28, right=1pt] {violated};
  \draw[aploop] (repair.west) -| (sim.south)
        node[aptagr, pos=0.72, left=1pt] {re-simulate};

  % ---------------- appearance joins the hand-off arrow from below ---------
  \draw[aparr] (texture.north) -- ($(render.east)!0.5!(gen.west)$);

  % ---------------- inline legend (below everything) -----------------------
  \coordinate (leg) at ($(spec.west |- repair.south)+(0,-0.62)$);
  \node[apsw, draw=adobeink, fill=adobegray!12, anchor=north west]
    (swA) at (leg) {};
  \node[apnote, right=3pt of swA] (lbA) {learned};

  \node[apsw, draw=adobeink, fill=white, right=9pt of lbA] (swB) {};
  \node[apnote, right=3pt of swB] (lbB) {deterministic};

  \node[apsw, draw=adobered, fill=white, right=9pt of lbB] (swC) {};
  \node[apnote, right=3pt of swC] (lbC) {trajectory check and re-author loop};

  \node[apsw, draw=adobegray, fill=white, dash pattern=on 2pt off 1.6pt,
        right=9pt of lbC] (swD) {};
  \node[apnote, right=3pt of swD] (lbD) {input};

\end{tikzpicture}
\endgroup

%       \caption{...}
%     \end{figure}
% The \label is carried HERE (see the block just below \makeatother), so main.tex
% must NOT add a second \label{fig:authoring}.
%
% CONTENT follows the subsection "From a physics description to a control video"
% (\ref{sec:gen-authoring}) and is drawn from the implementation in
% ~/diffphy_exp014/scripts/exp014/agentic_mask.py -- stage order from build_one()
% at lines 589-663, with:
%   text-in, no reference video   docstring l.4-8; author() l.294-306; main() l.683
%   scene specification written   author() l.294-327; scene contract l.124-149
%   required fields re-asked once _validate() l.249-265; author() l.322-326
%   waypoints only where needed   docstring l.26-29; motion model l.98-117
%   simulated by the renderer's   _sim_traj() l.348-352 (expand + simulate);
%     own engine                  silhouette_render.render_frames l.209 calls the
%                                 same simulate(); physics_world.py is pymunk
%   validation on the trajectory  lint_scene() l.355-450 (ballistic must be
%                                 dynamic l.406; off-frame l.410-417; rest /
%                                 support l.419-424; rising must rise l.426-428;
%                                 reach / miss targets l.430-445)
%   repair loop, three rounds     build_one() l.604-612; repair() l.472-488;
%                                 realized-motion readout _traj_summary l.453-469
%   appearance written apart      author_texture() l.272-291; called l.616
%   camera + graphic-body strip   apply_camera() l.491-504; _strip_graphic_bodies
%                                 l.210-215; called l.618-619
%   render + contact sheet        build_one l.623-627; render_gif; strip l.507-521
%   what the generator receives   build_one l.632-662 (fill-ratio cap l.632-638,
%                                 low-control-strength cap l.652-654), consumed by
%                                 generate_worker.py l.176-183 -> video_stylizer
%                                 WanVACEStylizer.stylize l.58-135
% The figure is descriptive of the method only: no accuracy, no comparison
% between configurations, and no quantity of any kind appears in it.
%
% The picture is a horizontal left-to-right pipeline about 12.9cm wide, so it
% stays inside the report's text width unscaled -- do not add scale= or
% \resizebox. Box labels are intentionally terse: the operational detail lives in
% the prose of \ref{sec:gen-authoring}, \ref{sec:gen-control} and
% \ref{sec:gen-mechanism}, so the figure carries only the dataflow -- the main
% spine (author -> simulate -> validate -> render -> generate), the bounded
% re-author loop dropping below the simulate/validate span, and the separate
% appearance track feeding the hand-off from below.
% =====================================================================
% The TikZ libraries must be loaded in the PREAMBLE, not here: a figure
% environment forms a group, and TikZ defines a library's macros locally, so a
% \usetikzlibrary issued inside a figure does not survive to the next one. This
% guard turns a missing preamble line into a clear message instead of TikZ's
% confusing "You need to say \usetikzlibrary{calc}" further down.
\makeatletter
\@ifundefined{tikz@library@calc@loaded}{%
  \PackageError{fig_authoring_pipeline}{%
    The preamble must contain\MessageBreak
    \string\usetikzlibrary{arrows.meta,positioning,calc}%
  }{Add that line after \string\usepackage{tikz}.}}{}
\makeatother

% The caller's \input runs BEFORE its \caption, so a bare \label here would bind
% to the enclosing section instead of to the figure. Step the figure counter,
% label that value, then hand the counter back: the \caption that follows steps
% it again to the same number, so \cref{fig:authoring} resolves to this figure.
% Outside a float the block does nothing.
\makeatletter
\@ifundefined{@captype}{}{%
  \refstepcounter{figure}%
  \label{fig:authoring}%
  \addtocounter{figure}{-1}%
}
\makeatother

\begingroup
\providecommand{\apsub}[1]{{\scriptsize\color{adobegray}#1}}

\begin{tikzpicture}[
    x=1cm, y=1cm,
    % --- shared box shape: terse title only; detail lives in the prose ----
    apbox/.style   = {rounded corners=2pt, line width=0.7pt, align=center,
                      inner xsep=6pt, inner ysep=4.2pt,
                      font=\sffamily\footnotesize, text=adobeink,
                      minimum height=8.5mm},
    % --- LEARNED component: light gray fill --------------------------------
    aplearn/.style = {apbox, draw=adobeink, fill=adobegray!12},
    % --- DETERMINISTIC component: white with an ink outline ----------------
    apdet/.style   = {apbox, draw=adobeink, fill=white},
    % --- accent: the check a trajectory must pass before it is rendered -----
    apgate/.style  = {apbox, draw=adobered, line width=0.9pt, fill=white,
                      text=adobered},
    % --- what is handed to the system (input) ------------------------------
    apdata/.style  = {apbox, draw=adobegray, fill=white,
                      dash pattern=on 2pt off 1.6pt},
    apline/.style  = {draw=adobeink, line width=0.7pt},
    aparr/.style   = {apline, -{Straight Barb[length=3.6pt,width=3.8pt]}},
    % the re-author loop, drawn in the accent colour so the cycle reads at a glance
    aploop/.style  = {draw=adobered, line width=0.7pt,
                      -{Straight Barb[length=3.6pt,width=3.8pt]}},
    apnote/.style  = {font=\sffamily\scriptsize, text=adobegray,
                      inner sep=0pt, align=left},
    aptag/.style   = {font=\sffamily\scriptsize, text=adobegray, fill=white,
                      inner xsep=3pt, inner ysep=1pt, align=center},
    aptagr/.style  = {aptag, text=adobered},
    apsw/.style    = {rounded corners=1.2pt, line width=0.7pt, inner sep=0pt,
                      minimum width=11pt, minimum height=7.5pt},
  ]

  % Layout: a horizontal main spine along y=0, left to right, sized to the text
  % width (five boxes: author -> simulate -> validate -> render -> generate).
  % The hand-off is folded into the render->generate arrow. The bounded
  % re-author loop drops below the simulate/validate span; the appearance track
  % sits below the render->generate span and joins the hand-off arrow from below.
  \def\bw{1.8cm}          % spine-box width
  \def\hg{0.46}           % horizontal gap between spine boxes
  \tikzset{spine/.style={apbox, text width=\bw, minimum height=13mm,
                         anchor=west}}

  % ---------------- main spine (left -> right) -----------------------------
  \node[spine, apdata] (spec)   at (0,0) {\textbf{Physics request}\\[1pt]\apsub{text}};
  \node[spine, aplearn] (author) at ($(spec.east)+(\hg,0)$)
    {\textbf{LM writes scene spec}};
  \node[spine, apdet] (sim) at ($(author.east)+(\hg,0)$)
    {\textbf{Simulate}\\[1pt]\apsub{renderer's engine}};
  \node[spine, apgate] (gate) at ($(sim.east)+(\hg,0)$)
    {\textbf{Validate trajectory}};
  \node[spine, apdet] (render) at ($(gate.east)+(\hg,0)$)
    {\textbf{Camera + render control}};
  \node[spine, aplearn] (gen) at ($(render.east)+(1.35,0)$)
    {\textbf{Diffusion backbone}};

  % ---------------- re-author loop (accent), below sim/validate ------------
  \node[aplearn, text width=3.4cm, anchor=north, minimum height=8.5mm]
    (repair) at ($(sim.south)!0.5!(gate.south)+(0,-0.92)$)
    {\textbf{Re-author}\ \ \apsub{$\leq 3$ rounds}};

  % ---------------- appearance track (learned), below render->gen ----------
  \node[aplearn, text width=3.7cm, anchor=north, minimum height=8.5mm]
    (texture) at ($(render.south east)!0.5!(gen.south west)+(0,-0.92)$)
    {\textbf{Appearance written separately}\\[1pt]\apsub{no access to the control}};

  % ---------------- edges: main spine --------------------------------------
  \draw[aparr] (spec)    -- (author);
  \draw[aparr] (author)  -- (sim);
  \draw[aparr] (sim)     -- (gate);
  \draw[aparr] (gate)    -- (render)
        node[aptagr, midway, above=0pt] {validated};
  % render -> generate carries the hand-off (control + appearance + seed + strength)
  \draw[aparr] (render)  -- (gen)
        node[aptag, midway, above=0pt] {hand off};

  % ---------------- re-author loop: validate -> re-author -> re-simulate ----
  \draw[aploop] (gate.south) |- (repair.east)
        node[aptagr, pos=0.28, right=1pt] {violated};
  \draw[aploop] (repair.west) -| (sim.south)
        node[aptagr, pos=0.72, left=1pt] {re-simulate};

  % ---------------- appearance joins the hand-off arrow from below ---------
  \draw[aparr] (texture.north) -- ($(render.east)!0.5!(gen.west)$);

  % ---------------- inline legend (below everything) -----------------------
  \coordinate (leg) at ($(spec.west |- repair.south)+(0,-0.62)$);
  \node[apsw, draw=adobeink, fill=adobegray!12, anchor=north west]
    (swA) at (leg) {};
  \node[apnote, right=3pt of swA] (lbA) {learned};

  \node[apsw, draw=adobeink, fill=white, right=9pt of lbA] (swB) {};
  \node[apnote, right=3pt of swB] (lbB) {deterministic};

  \node[apsw, draw=adobered, fill=white, right=9pt of lbB] (swC) {};
  \node[apnote, right=3pt of swC] (lbC) {trajectory check and re-author loop};

  \node[apsw, draw=adobegray, fill=white, dash pattern=on 2pt off 1.6pt,
        right=9pt of lbC] (swD) {};
  \node[apnote, right=3pt of swD] (lbD) {input};

\end{tikzpicture}
\endgroup

%       \caption{...}
%     \end{figure}
% The \label is carried HERE (see the block just below \makeatother), so main.tex
% must NOT add a second \label{fig:authoring}.
%
% CONTENT follows the subsection "From a physics description to a control video"
% (\ref{sec:gen-authoring}) and is drawn from the implementation in
% ~/diffphy_exp014/scripts/exp014/agentic_mask.py -- stage order from build_one()
% at lines 589-663, with:
%   text-in, no reference video   docstring l.4-8; author() l.294-306; main() l.683
%   scene specification written   author() l.294-327; scene contract l.124-149
%   required fields re-asked once _validate() l.249-265; author() l.322-326
%   waypoints only where needed   docstring l.26-29; motion model l.98-117
%   simulated by the renderer's   _sim_traj() l.348-352 (expand + simulate);
%     own engine                  silhouette_render.render_frames l.209 calls the
%                                 same simulate(); physics_world.py is pymunk
%   validation on the trajectory  lint_scene() l.355-450 (ballistic must be
%                                 dynamic l.406; off-frame l.410-417; rest /
%                                 support l.419-424; rising must rise l.426-428;
%                                 reach / miss targets l.430-445)
%   repair loop, three rounds     build_one() l.604-612; repair() l.472-488;
%                                 realized-motion readout _traj_summary l.453-469
%   appearance written apart      author_texture() l.272-291; called l.616
%   camera + graphic-body strip   apply_camera() l.491-504; _strip_graphic_bodies
%                                 l.210-215; called l.618-619
%   render + contact sheet        build_one l.623-627; render_gif; strip l.507-521
%   what the generator receives   build_one l.632-662 (fill-ratio cap l.632-638,
%                                 low-control-strength cap l.652-654), consumed by
%                                 generate_worker.py l.176-183 -> video_stylizer
%                                 WanVACEStylizer.stylize l.58-135
% The figure is descriptive of the method only: no accuracy, no comparison
% between configurations, and no quantity of any kind appears in it.
%
% The picture is a horizontal left-to-right pipeline about 12.9cm wide, so it
% stays inside the report's text width unscaled -- do not add scale= or
% \resizebox. Box labels are intentionally terse: the operational detail lives in
% the prose of \ref{sec:gen-authoring}, \ref{sec:gen-control} and
% \ref{sec:gen-mechanism}, so the figure carries only the dataflow -- the main
% spine (author -> simulate -> validate -> render -> generate), the bounded
% re-author loop dropping below the simulate/validate span, and the separate
% appearance track feeding the hand-off from below.
% =====================================================================
% The TikZ libraries must be loaded in the PREAMBLE, not here: a figure
% environment forms a group, and TikZ defines a library's macros locally, so a
% \usetikzlibrary issued inside a figure does not survive to the next one. This
% guard turns a missing preamble line into a clear message instead of TikZ's
% confusing "You need to say \usetikzlibrary{calc}" further down.
\makeatletter
\@ifundefined{tikz@library@calc@loaded}{%
  \PackageError{fig_authoring_pipeline}{%
    The preamble must contain\MessageBreak
    \string\usetikzlibrary{arrows.meta,positioning,calc}%
  }{Add that line after \string\usepackage{tikz}.}}{}
\makeatother

% The caller's \input runs BEFORE its \caption, so a bare \label here would bind
% to the enclosing section instead of to the figure. Step the figure counter,
% label that value, then hand the counter back: the \caption that follows steps
% it again to the same number, so \cref{fig:authoring} resolves to this figure.
% Outside a float the block does nothing.
\makeatletter
\@ifundefined{@captype}{}{%
  \refstepcounter{figure}%
  \label{fig:authoring}%
  \addtocounter{figure}{-1}%
}
\makeatother

\begingroup
\providecommand{\apsub}[1]{{\scriptsize\color{adobegray}#1}}

\begin{tikzpicture}[
    x=1cm, y=1cm,
    % --- shared box shape: terse title only; detail lives in the prose ----
    apbox/.style   = {rounded corners=2pt, line width=0.7pt, align=center,
                      inner xsep=6pt, inner ysep=4.2pt,
                      font=\sffamily\footnotesize, text=adobeink,
                      minimum height=8.5mm},
    % --- LEARNED component: light gray fill --------------------------------
    aplearn/.style = {apbox, draw=adobeink, fill=adobegray!12},
    % --- DETERMINISTIC component: white with an ink outline ----------------
    apdet/.style   = {apbox, draw=adobeink, fill=white},
    % --- accent: the check a trajectory must pass before it is rendered -----
    apgate/.style  = {apbox, draw=adobered, line width=0.9pt, fill=white,
                      text=adobered},
    % --- what is handed to the system (input) ------------------------------
    apdata/.style  = {apbox, draw=adobegray, fill=white,
                      dash pattern=on 2pt off 1.6pt},
    apline/.style  = {draw=adobeink, line width=0.7pt},
    aparr/.style   = {apline, -{Straight Barb[length=3.6pt,width=3.8pt]}},
    % the re-author loop, drawn in the accent colour so the cycle reads at a glance
    aploop/.style  = {draw=adobered, line width=0.7pt,
                      -{Straight Barb[length=3.6pt,width=3.8pt]}},
    apnote/.style  = {font=\sffamily\scriptsize, text=adobegray,
                      inner sep=0pt, align=left},
    aptag/.style   = {font=\sffamily\scriptsize, text=adobegray, fill=white,
                      inner xsep=3pt, inner ysep=1pt, align=center},
    aptagr/.style  = {aptag, text=adobered},
    apsw/.style    = {rounded corners=1.2pt, line width=0.7pt, inner sep=0pt,
                      minimum width=11pt, minimum height=7.5pt},
  ]

  % Layout: a horizontal main spine along y=0, left to right, sized to the text
  % width (five boxes: author -> simulate -> validate -> render -> generate).
  % The hand-off is folded into the render->generate arrow. The bounded
  % re-author loop drops below the simulate/validate span; the appearance track
  % sits below the render->generate span and joins the hand-off arrow from below.
  \def\bw{1.8cm}          % spine-box width
  \def\hg{0.46}           % horizontal gap between spine boxes
  \tikzset{spine/.style={apbox, text width=\bw, minimum height=13mm,
                         anchor=west}}

  % ---------------- main spine (left -> right) -----------------------------
  \node[spine, apdata] (spec)   at (0,0) {\textbf{Physics request}\\[1pt]\apsub{text}};
  \node[spine, aplearn] (author) at ($(spec.east)+(\hg,0)$)
    {\textbf{LM writes scene spec}};
  \node[spine, apdet] (sim) at ($(author.east)+(\hg,0)$)
    {\textbf{Simulate}\\[1pt]\apsub{renderer's engine}};
  \node[spine, apgate] (gate) at ($(sim.east)+(\hg,0)$)
    {\textbf{Validate trajectory}};
  \node[spine, apdet] (render) at ($(gate.east)+(\hg,0)$)
    {\textbf{Camera + render control}};
  \node[spine, aplearn] (gen) at ($(render.east)+(1.35,0)$)
    {\textbf{Diffusion backbone}};

  % ---------------- re-author loop (accent), below sim/validate ------------
  \node[aplearn, text width=3.4cm, anchor=north, minimum height=8.5mm]
    (repair) at ($(sim.south)!0.5!(gate.south)+(0,-0.92)$)
    {\textbf{Re-author}\ \ \apsub{$\leq 3$ rounds}};

  % ---------------- appearance track (learned), below render->gen ----------
  \node[aplearn, text width=3.7cm, anchor=north, minimum height=8.5mm]
    (texture) at ($(render.south east)!0.5!(gen.south west)+(0,-0.92)$)
    {\textbf{Appearance written separately}\\[1pt]\apsub{no access to the control}};

  % ---------------- edges: main spine --------------------------------------
  \draw[aparr] (spec)    -- (author);
  \draw[aparr] (author)  -- (sim);
  \draw[aparr] (sim)     -- (gate);
  \draw[aparr] (gate)    -- (render)
        node[aptagr, midway, above=0pt] {validated};
  % render -> generate carries the hand-off (control + appearance + seed + strength)
  \draw[aparr] (render)  -- (gen)
        node[aptag, midway, above=0pt] {hand off};

  % ---------------- re-author loop: validate -> re-author -> re-simulate ----
  \draw[aploop] (gate.south) |- (repair.east)
        node[aptagr, pos=0.28, right=1pt] {violated};
  \draw[aploop] (repair.west) -| (sim.south)
        node[aptagr, pos=0.72, left=1pt] {re-simulate};

  % ---------------- appearance joins the hand-off arrow from below ---------
  \draw[aparr] (texture.north) -- ($(render.east)!0.5!(gen.west)$);

  % ---------------- inline legend (below everything) -----------------------
  \coordinate (leg) at ($(spec.west |- repair.south)+(0,-0.62)$);
  \node[apsw, draw=adobeink, fill=adobegray!12, anchor=north west]
    (swA) at (leg) {};
  \node[apnote, right=3pt of swA] (lbA) {learned};

  \node[apsw, draw=adobeink, fill=white, right=9pt of lbA] (swB) {};
  \node[apnote, right=3pt of swB] (lbB) {deterministic};

  \node[apsw, draw=adobered, fill=white, right=9pt of lbB] (swC) {};
  \node[apnote, right=3pt of swC] (lbC) {trajectory check and re-author loop};

  \node[apsw, draw=adobegray, fill=white, dash pattern=on 2pt off 1.6pt,
        right=9pt of lbC] (swD) {};
  \node[apnote, right=3pt of swD] (lbD) {input};

\end{tikzpicture}
\endgroup

%       \caption{...}
%     \end{figure}
% The \label is carried HERE (see the block just below \makeatother), so main.tex
% must NOT add a second \label{fig:authoring}.
%
% CONTENT follows the subsection "From a physics description to a control video"
% (\ref{sec:gen-authoring}) and is drawn from the implementation in
% ~/diffphy_exp014/scripts/exp014/agentic_mask.py -- stage order from build_one()
% at lines 589-663, with:
%   text-in, no reference video   docstring l.4-8; author() l.294-306; main() l.683
%   scene specification written   author() l.294-327; scene contract l.124-149
%   required fields re-asked once _validate() l.249-265; author() l.322-326
%   waypoints only where needed   docstring l.26-29; motion model l.98-117
%   simulated by the renderer's   _sim_traj() l.348-352 (expand + simulate);
%     own engine                  silhouette_render.render_frames l.209 calls the
%                                 same simulate(); physics_world.py is pymunk
%   validation on the trajectory  lint_scene() l.355-450 (ballistic must be
%                                 dynamic l.406; off-frame l.410-417; rest /
%                                 support l.419-424; rising must rise l.426-428;
%                                 reach / miss targets l.430-445)
%   repair loop, three rounds     build_one() l.604-612; repair() l.472-488;
%                                 realized-motion readout _traj_summary l.453-469
%   appearance written apart      author_texture() l.272-291; called l.616
%   camera + graphic-body strip   apply_camera() l.491-504; _strip_graphic_bodies
%                                 l.210-215; called l.618-619
%   render + contact sheet        build_one l.623-627; render_gif; strip l.507-521
%   what the generator receives   build_one l.632-662 (fill-ratio cap l.632-638,
%                                 low-control-strength cap l.652-654), consumed by
%                                 generate_worker.py l.176-183 -> video_stylizer
%                                 WanVACEStylizer.stylize l.58-135
% The figure is descriptive of the method only: no accuracy, no comparison
% between configurations, and no quantity of any kind appears in it.
%
% The picture is a horizontal left-to-right pipeline about 12.9cm wide, so it
% stays inside the report's text width unscaled -- do not add scale= or
% \resizebox. Box labels are intentionally terse: the operational detail lives in
% the prose of \ref{sec:gen-authoring}, \ref{sec:gen-control} and
% \ref{sec:gen-mechanism}, so the figure carries only the dataflow -- the main
% spine (author -> simulate -> validate -> render -> generate), the bounded
% re-author loop dropping below the simulate/validate span, and the separate
% appearance track feeding the hand-off from below.
% =====================================================================
% The TikZ libraries must be loaded in the PREAMBLE, not here: a figure
% environment forms a group, and TikZ defines a library's macros locally, so a
% \usetikzlibrary issued inside a figure does not survive to the next one. This
% guard turns a missing preamble line into a clear message instead of TikZ's
% confusing "You need to say \usetikzlibrary{calc}" further down.
\makeatletter
\@ifundefined{tikz@library@calc@loaded}{%
  \PackageError{fig_authoring_pipeline}{%
    The preamble must contain\MessageBreak
    \string\usetikzlibrary{arrows.meta,positioning,calc}%
  }{Add that line after \string\usepackage{tikz}.}}{}
\makeatother

% The caller's \input runs BEFORE its \caption, so a bare \label here would bind
% to the enclosing section instead of to the figure. Step the figure counter,
% label that value, then hand the counter back: the \caption that follows steps
% it again to the same number, so \cref{fig:authoring} resolves to this figure.
% Outside a float the block does nothing.
\makeatletter
\@ifundefined{@captype}{}{%
  \refstepcounter{figure}%
  \label{fig:authoring}%
  \addtocounter{figure}{-1}%
}
\makeatother

\begingroup
\providecommand{\apsub}[1]{{\scriptsize\color{adobegray}#1}}

\begin{tikzpicture}[
    x=1cm, y=1cm,
    % --- shared box shape: terse title only; detail lives in the prose ----
    apbox/.style   = {rounded corners=2pt, line width=0.7pt, align=center,
                      inner xsep=6pt, inner ysep=4.2pt,
                      font=\sffamily\footnotesize, text=adobeink,
                      minimum height=8.5mm},
    % --- LEARNED component: light gray fill --------------------------------
    aplearn/.style = {apbox, draw=adobeink, fill=adobegray!12},
    % --- DETERMINISTIC component: white with an ink outline ----------------
    apdet/.style   = {apbox, draw=adobeink, fill=white},
    % --- accent: the check a trajectory must pass before it is rendered -----
    apgate/.style  = {apbox, draw=adobered, line width=0.9pt, fill=white,
                      text=adobered},
    % --- what is handed to the system (input) ------------------------------
    apdata/.style  = {apbox, draw=adobegray, fill=white,
                      dash pattern=on 2pt off 1.6pt},
    apline/.style  = {draw=adobeink, line width=0.7pt},
    aparr/.style   = {apline, -{Straight Barb[length=3.6pt,width=3.8pt]}},
    % the re-author loop, drawn in the accent colour so the cycle reads at a glance
    aploop/.style  = {draw=adobered, line width=0.7pt,
                      -{Straight Barb[length=3.6pt,width=3.8pt]}},
    apnote/.style  = {font=\sffamily\scriptsize, text=adobegray,
                      inner sep=0pt, align=left},
    aptag/.style   = {font=\sffamily\scriptsize, text=adobegray, fill=white,
                      inner xsep=3pt, inner ysep=1pt, align=center},
    aptagr/.style  = {aptag, text=adobered},
    apsw/.style    = {rounded corners=1.2pt, line width=0.7pt, inner sep=0pt,
                      minimum width=11pt, minimum height=7.5pt},
  ]

  % Layout: a horizontal main spine along y=0, left to right, sized to the text
  % width (five boxes: author -> simulate -> validate -> render -> generate).
  % The hand-off is folded into the render->generate arrow. The bounded
  % re-author loop drops below the simulate/validate span; the appearance track
  % sits below the render->generate span and joins the hand-off arrow from below.
  \def\bw{1.8cm}          % spine-box width
  \def\hg{0.46}           % horizontal gap between spine boxes
  \tikzset{spine/.style={apbox, text width=\bw, minimum height=13mm,
                         anchor=west}}

  % ---------------- main spine (left -> right) -----------------------------
  \node[spine, apdata] (spec)   at (0,0) {\textbf{Physics request}\\[1pt]\apsub{text}};
  \node[spine, aplearn] (author) at ($(spec.east)+(\hg,0)$)
    {\textbf{LM writes scene spec}};
  \node[spine, apdet] (sim) at ($(author.east)+(\hg,0)$)
    {\textbf{Simulate}\\[1pt]\apsub{renderer's engine}};
  \node[spine, apgate] (gate) at ($(sim.east)+(\hg,0)$)
    {\textbf{Validate trajectory}};
  \node[spine, apdet] (render) at ($(gate.east)+(\hg,0)$)
    {\textbf{Camera + render control}};
  \node[spine, aplearn] (gen) at ($(render.east)+(1.35,0)$)
    {\textbf{Diffusion backbone}};

  % ---------------- re-author loop (accent), below sim/validate ------------
  \node[aplearn, text width=3.4cm, anchor=north, minimum height=8.5mm]
    (repair) at ($(sim.south)!0.5!(gate.south)+(0,-0.92)$)
    {\textbf{Re-author}\ \ \apsub{$\leq 3$ rounds}};

  % ---------------- appearance track (learned), below render->gen ----------
  \node[aplearn, text width=3.7cm, anchor=north, minimum height=8.5mm]
    (texture) at ($(render.south east)!0.5!(gen.south west)+(0,-0.92)$)
    {\textbf{Appearance written separately}\\[1pt]\apsub{no access to the control}};

  % ---------------- edges: main spine --------------------------------------
  \draw[aparr] (spec)    -- (author);
  \draw[aparr] (author)  -- (sim);
  \draw[aparr] (sim)     -- (gate);
  \draw[aparr] (gate)    -- (render)
        node[aptagr, midway, above=0pt] {validated};
  % render -> generate carries the hand-off (control + appearance + seed + strength)
  \draw[aparr] (render)  -- (gen)
        node[aptag, midway, above=0pt] {hand off};

  % ---------------- re-author loop: validate -> re-author -> re-simulate ----
  \draw[aploop] (gate.south) |- (repair.east)
        node[aptagr, pos=0.28, right=1pt] {violated};
  \draw[aploop] (repair.west) -| (sim.south)
        node[aptagr, pos=0.72, left=1pt] {re-simulate};

  % ---------------- appearance joins the hand-off arrow from below ---------
  \draw[aparr] (texture.north) -- ($(render.east)!0.5!(gen.west)$);

  % ---------------- inline legend (below everything) -----------------------
  \coordinate (leg) at ($(spec.west |- repair.south)+(0,-0.62)$);
  \node[apsw, draw=adobeink, fill=adobegray!12, anchor=north west]
    (swA) at (leg) {};
  \node[apnote, right=3pt of swA] (lbA) {learned};

  \node[apsw, draw=adobeink, fill=white, right=9pt of lbA] (swB) {};
  \node[apnote, right=3pt of swB] (lbB) {deterministic};

  \node[apsw, draw=adobered, fill=white, right=9pt of lbB] (swC) {};
  \node[apnote, right=3pt of swC] (lbC) {trajectory check and re-author loop};

  \node[apsw, draw=adobegray, fill=white, dash pattern=on 2pt off 1.6pt,
        right=9pt of lbC] (swD) {};
  \node[apnote, right=3pt of swD] (lbD) {input};

\end{tikzpicture}
\endgroup
